% ============================================================
%  Section 2 – Background
% ============================================================
\subsection{GGM trees}
The Goldreich--Goldwasser--Micali construction~\cite{GGM86} builds a
pseudorandom function from a pseudorandom generator by iterated expansion along
the branches of a binary tree indexed by the input bits. It is the construction
studied throughout this work.


\begin{definition}[GGM Tree]
\label{def:ggm}
Let $\PRG : \{0,1\}^{\lam} \to \{0,1\}^{2\lam}$ be a length-doubling pseudorandom
generator, decomposed as $\PRG(x) = (\PRGl(x),\, \PRGr(x))$, where $\PRGl$ and
$\PRGr$ each output $\lam$ bits.  Given a root seed
$\seed_r \in \{0,1\}^{\lam}$ and a depth $d$, the \emph{GGM tree} is a complete
binary tree of depth $d$ in which:
\begin{itemize}
  \item the root holds $\seed_r$;
  \item each internal node $v$ storing seed $s$ generates its left child as
        $\PRGl(s)$ and its right child as $\PRGr(s)$;
  \item the $2^{d}$ leaves are the \PRG\ outputs at depth $d$.
\end{itemize}
\end{definition}

\revon
For a secure PRG, the seed of a punctured leaf remains computationally
hidden even when the legitimate seeds of other leaves are disclosed.
Structurally, a published ancestor permits forward derivation of every leaf
in its subtree. These computational and structural statements are distinct
(see Figure~\ref{fig:less-trees}).
\revoff

\paragraph{Incomplete \GGM\ trees.}
When the number of required leaves $t$ is not a power of $2$, a scheme may use an
\emph{incomplete} tree: leaves appear at multiple levels, and internal nodes near
the bottom may have fewer than two children or act as leaves themselves. This
reduces memory and computation, but the additional index-tracking logic introduces
new fault surfaces. The scheme may also use the first $t$ leaves of the tree.
%as detailed in Section~\ref{sec:theoretical}.

% ============================================================
\paragraph{The common template.}
\label{sec:template}
% ============================================================

All schemes considered here share one skeleton, worth stating once. A $3$- or
$5$-pass identification protocol is repeated $t$ times in parallel and made
non-interactive by Fiat--Shamir. Round $i$ needs a fresh ephemeral secret,
derived from a leaf seed $\eseed_i$; the $t$ leaf seeds are the leaves of one
\GGM\ tree, so the signer stores $\lambda$ bits instead of $t\lambda$.

The challenge $\mathbf{c}$ has fixed weight $w$ and splits the rounds in two.
For the $w$ \emph{hidden} rounds ($i \in \calH$) the signer publishes a
response $\resp_i$ combining the ephemeral secret with the long-term key, so
$\eseed_i$ must stay secret: $\resp_i$ and $\eseed_i$ together determine that
key. For the $t-w$ \emph{revealed} rounds ($i \in \calR$) the seed itself is
released and the verifier recomputes the commitment.

Releasing $t-w$ individual seeds would undo the compression, so the signer
publishes instead the boundary of the hidden subtree --- a small set of
internal nodes from which exactly the revealed leaves, and no others, can be
recomputed. We call this set the \emph{path}. Computing it is the only place
in signing where the code decides, per node, whether a seed may be released,
and that decision is the object of this paper: \chg{it is what the attacks
of~\cite{ZKFault,FaultToForge2025}, and every attack we mount in
Section~\ref{sec:app}, corrupt. The remaining attacks surveyed in
Section~\ref{sec:relwork} reach the same leak through other components;
Section~\ref{sec:taxonomy} places them.}

% ============================================================
\subsection{The MEDS signature scheme}
\label{sec:meds}
% ============================================================

The Matrix Equivalence Digital Signature (MEDS) scheme~\cite{ChouNPRRST23} is
a code-based signature whose security rests on the hardness of the Matrix Code
Equivalence (MCE) problem. Because the practical attacks of
Section~\ref{sec:app} target the MEDS reference implementation, we describe it
in more detail than the other instances.

\paragraph{Matrix codes and the MCE problem.}
A \emph{matrix code} of dimension $k$ is a $k$-dimensional $\Fq$-subspace
$\Mat \subseteq \Fq^{m \times n}$. Two matrix codes $\Mat, \Matd$ are
\emph{equivalent} if there exist $A \in \GL_m(q)$ and $B \in \GL_n(q)$ with
\[
  \Matd \;=\; A\,\Mat\,B \;:=\; \{\, A\,C\,B \;:\; C \in \Mat \,\}.
\]

\begin{definition}[Matrix Code Equivalence]
\label{def:mce}
Given two $k$-dimensional matrix codes $\Mat, \Matd \subseteq \Fq^{m\times n}$,
the \emph{MCE problem} asks for a pair $(A,B) \in \GL_m(q) \times \GL_n(q)$
such that $\Matd = A\,\Mat\,B$.
\end{definition}

In implementations a matrix code is handled through a generator matrix
$G \in \Fq^{k \times mn}$ whose $k$ rows are the vectorisations
$\vect(C_1),\ldots,\vect(C_k)$ of a basis of $\Mat$. Write $\pimap{A}{B}$ for
the induced row-wise map
\[
  \pimap{A}{B}\bigl(\vect(C)\bigr) \;=\; \vect\bigl(A\,C\,B\bigr),
\]
extended to generator matrices by applying it to each row. Since a generator
matrix fixes a basis but equivalence only relates the underlying
\emph{subspaces}, the relation between two generator matrices of equivalent
codes carries a third, unknown ingredient: a change of basis
$T \in \GL_k(q)$. Concretely, $G_0$ and $G_1$ generate equivalent codes if and
only if there exists a triple $(A,B,T)$ with
\begin{equation}
  \label{eq:mce-gen}
  G_1 \;=\; T \cdot \pimap{A}{B}(G_0),
  \qquad\text{equivalently}\qquad
  G_1 \;=\; \SF\bigl(\pimap{A}{B}(G_0)\bigr),
\end{equation}
where $\SF$ denotes the systematic form, which absorbs $T$ by fixing a
canonical basis.

\revon
\chg{The unknown basis change $T$ is central to the hardness: with $T$ fixed the
\revoff
relations are linear in $(A,B^{-1})$, which MEDS key generation exploits to
sample keys, whereas with $T$ unknown \eqref{eq:mce-gen} is a system of
bilinear equations (Remark~\ref{rem:mce-hardness} in
Appendix~\ref{app:algs}). MCE is Tensor-Isomorphism-complete~\cite{GrochowQiao23};
see~\cite{ReijndersST23,NarayananQT24} for the best known classical
algorithms. Subroutine names follow the reference
specification~\cite{ChouNPRRST23,MEDS_PQC} and are summarised in
Table~\ref{tab:meds-notation} (Appendix~\ref{app:algs}).}

\paragraph{Key generation.}
Key generation (Algorithm~\ref{alg:keygen_meds}, Appendix~\ref{app:algs})
expands a public seed into $G_0$ and then produces $s-1$ equivalent codes: for
each $i$ it samples $T_i \in \GL_k(q)$, recovers $(A_i,B_i^{-1})$ from
$T_i G_0$ by the linear solve of Remark~\ref{rem:mce-hardness}, and publishes
$G_i \leftarrow \SF(\pimap{A_i}{B_i}(G_0))$. The secret key retains the pairs
\revon
$(A_i^{-1},B_i^{-1})$; these pairs are sufficient to construct signing-equivalent secret information.
\revoff

\paragraph{Signing and the role of the seed tree.}
Algorithm~\ref{alg:sign_meds} follows the template just described. A fresh $\delta$ is expanded into a root seed
$\rho$ and a salt $\alpha$; \chg{line~\ref{alg:meds:tree}} grows the \GGM\ tree of
Definition~\ref{def:ggm} from $\rho$, yielding the $t$ ephemeral leaf seeds
$\sigma_0,\ldots,\sigma_{t-1}$. Each $\sigma_i$ is expanded into an ephemeral
equivalence $(\tilde{A}_i, \tilde{B}_i)$, and round $i$ commits to the
resulting generator matrix
$\tilde{G}_i = \SF(\pimap{\tilde{A}_i}{\tilde{B}_i}(G_0))$ by hashing its
non-systematic part. Hashing all $t$ commitments together with the message
gives the digest $d$, parsed into the weight-$w$ challenge
$\mathbf{h} = (h_0,\ldots,h_{t-1})$.

For each hidden round ($h_i > 0$) the signer publishes the response pair
\begin{equation}
  \label{eq:meds-resp}
  (\mu_i, \nu_i) \;=\; \bigl(\tilde{A}_i A_{h_i}^{-1},\;
                             B_{h_i}^{-1} \tilde{B}_i\bigr),
\end{equation}
and for the revealed rounds the seeds are released compactly through the path
$p \leftarrow \SeedTreeToPath_t(\mathbf{h},\rho,\alpha)$ computed at
\chg{line~\ref{alg:meds:path}},
the target of every attack in Section~\ref{sec:app}.

\begin{algorithm}[tb]
\footnotesize
\setlength{\algomargin}{0.9em}
\caption{MEDS signing (simplified).}
\label{alg:sign_meds}
\KwIn{Secret key $\sk$, message $m$}
\KwOut{Signed message $m_s$}

Parse $\sk$ to recover $\sigma_{G_0}$, $A_i^{-1}$, $B_i^{-1}$, $i=1,\ldots,s-1$;\quad
$G_0 \leftarrow \ExpandSystMat(\sigma_{G_0})$\;
$\delta \xleftarrow{\$} \mathcal{B}^{\ell_{\mathrm{sec\_seed}}}$;\quad
$\rho,\, \alpha \leftarrow \XOF(\delta,\, \ell_{\mathrm{tree\_seed}},\, \ell_{\mathrm{salt}})$\;
$(\sigma_0, \ldots, \sigma_{t-1}) \leftarrow \SeedTree_t(\rho,\, \alpha)$\label{alg:meds:tree}
\tcp*[r]{\GGM\ tree, Def.~\ref{def:ggm}}

\For{$i = 0$ \KwTo $t-1$}{
    $\sigma'_i \leftarrow \bigl(\alpha \mid \sigma_i \mid \ToBytes(2^{1+\lceil\log_2 t\rceil}+i,4)\bigr)$;\quad
    $\sigma_{\tilde{A}_i},\, \sigma_{\tilde{B}_i} \leftarrow \XOF(\sigma'_i,\, \ell_{\mathrm{pub\_seed}},\, \ell_{\mathrm{pub\_seed}})$\;
    $\tilde{A}_i \leftarrow \ExpandInvMat(\sigma_{\tilde{A}_i}, m)$;\quad
    $\tilde{B}_i \leftarrow \ExpandInvMat(\sigma_{\tilde{B}_i}, n)$;\quad
    $\tilde{G}_i \leftarrow \SF\!\bigl(\pimap{\tilde{A}_i}{\tilde{B}_i}(G_0)\bigr)$
    \tcp*[r]{$\perp\Rightarrow$ resample}
}

$d \leftarrow H\!\bigl(\Compress(\tilde{G}_0[;k,mn{-}1]) \mid \cdots \mid \Compress(\tilde{G}_{t-1}[;k,mn{-}1]) \mid m\bigr)$\;
$(h_0,\ldots,h_{t-1}) \leftarrow \ParseHash_{s,t,w}(d)$
\tcp*[r]{weight-$w$ challenge}

\For{$i = 0$ \KwTo $t-1$ \textbf{with} $h_i > 0$}{
    $\mu_i \leftarrow \tilde{A}_i A_{h_i}^{-1}$;\quad
    $\nu_i \leftarrow B_{h_i}^{-1} \tilde{B}_i$;\quad
    append $\bigl(\Compress(\mu_i) \mid \Compress(\nu_i)\bigr)$ to $v$\;
}

$p \leftarrow \SeedTreeToPath_t(h_0,\ldots,h_{t-1},\rho,\alpha)$\label{alg:meds:path}
\tcp*[r]{\textbf{attack target}}
\Return $m_s \leftarrow \bigl(v \mid p \mid d \mid \alpha \mid m\bigr)$
\end{algorithm}

\paragraph{Verification.}
Verification (Algorithm~\ref{alg:verify_meds}, Appendix~\ref{app:algs})
recomputes the challenge from $d$, runs $\PathToSeedTree_t$ to recover the
$t-w$ revealed seeds, rebuilds every $\hat{G}_i$---from $(\mu_i,\nu_i)$ and
$G_{h_i}$ when $h_i > 0$, from the recovered seed when $h_i = 0$---and accepts
iff the recomputed digest equals $d$. Note that for hidden rounds the verifier
never uses $\sigma_i$: the seed of a hidden round is simply absent from an
honest signature.

\paragraph{Secret-extraction opportunity.}
This last observation is what the attacks exploit. Fix a hidden round $i$ with
$h_i = j$. The signature already contains $(\mu_i,\nu_i)$
from~\eqref{eq:meds-resp}. If an attacker additionally obtains $\sigma_i$, it
expands it into $(\tilde{A}_i,\tilde{B}_i)$ and reads off
\begin{equation}
  \label{eq:meds-extract}
  A_j^{-1} \;=\; \tilde{A}_i^{-1} \mu_i,
  \qquad
  B_j^{-1} \;=\; \nu_i \tilde{B}_i^{-1},
\end{equation}
\revon
by two matrix inversions and two multiplications over $\Fq$. Correct seed publication prevents this pair of revelations during
an honest signing execution.
\revoff

\begin{table}[htbp]
\centering
\caption{MEDS parameter sets~\cite{ChouNPRRST23}; MEDS uses $n{=}m{=}k$. \chg{For
an effective root fault, the last three columns give
$p_{\min} = 1-((s-2)/(s-1))^{w}$ of Proposition~\ref{prop:alg-correctness},
\revon
the probability $\Pr[Q{=}1]$ that one \emph{effective root disclosure}
covers every long-term matrix pair, and the exact expected number $\mathbb{E}[Q]$ of effective faulted
\revoff
signatures needed (Appendix~\ref{app:proof-query}).}}
\label{tab:scheme-params}
\setlength{\tabcolsep}{4pt}
\small
\begin{tabular}{@{}lrrrcl rrr@{}}
\toprule
\textbf{Parameter set} & $t$ & $w$ & $s$ & \chg{\textbf{Level}} & \textbf{Field}
 & $p_{\min}$ & \chg{$\Pr[Q{=}1]$} & \chg{$\mathbb{E}[Q]$} \\
\midrule
MEDS-9923   & 1152 & 14 & 4 & I   & $n{=}14,\,q{=}4093$ & 0.9966 & \chg{0.9897} & \chg{1.0103} \\
MEDS-13220  &  192 & 20 & 5 & I   & $n{=}14,\,q{=}4093$ & 0.9968 & \chg{0.9873} & \chg{1.0127} \\
MEDS-41711  &  608 & 26 & 4 & III & $n{=}22,\,q{=}4093$ & 1.0000 & \chg{0.9999} & \chg{1.0001} \\
MEDS-69497  &  160 & 36 & 5 & III & $n{=}22,\,q{=}4093$ & 1.0000 & \chg{0.9999} & \chg{1.0001} \\
MEDS-134180 &  192 & 52 & 5 & V   & $n{=}30,\,q{=}2039$ & 1.0000 & \chg{1.0000} & \chg{1.0000} \\
MEDS-167717 &  112 & 66 & 6 & V   & $n{=}30,\,q{=}2039$ & 1.0000 & \chg{1.0000} & \chg{1.0000} \\
\bottomrule
\end{tabular}
\end{table}

% ============================================================
\subsection{Other instances: LESS-v2, CROSS, and MPCitH}
\label{sec:others}
% ============================================================

The other schemes differ in the algebra but not in the mechanism, and we
record only what the fault analysis needs.

\textbf{LESS-v2}~\cite{LESS-spec} applies Fiat--Shamir to a $3$-pass sigma
protocol over the Canonical Form Linear Equivalence problem. Its long-term
secret is a set of monomial matrices $\pi_1,\ldots,\pi_{s-1} \in \Mon_n(q)$,
and a hidden round publishes the coset representative
$\CosetRep(\pi_{b[i]} \circ \tilde{\pi}^*_i)$. Two structural differences from
MEDS matter, because the attacks of~\cite{ZKFault,FaultToForge2025} turn on
them: the tree is \emph{incomplete}, so leaves sit at two depths and the
implementation carries \chg{per-level leaf-offset arrays} to locate them; and the
path is computed from an explicit auxiliary \emph{flag tree}, with
$\LabelLeaves$ writing $\calF(\mathrm{leaf}_i) = \mathbf{1}[i \in \calH]$ and
$\ComputeSeeds$ propagating $\calF(v) = \calF(\ell) \vee \calF(r)$ upwards.
Both the \chg{leaf-offset arrays} and the flag array are state that MEDS does not
have, and hence fault surface that MEDS does not have\chg{; MEDS has a different
piece of state instead, a cursor into the sorted list of hidden leaves
(Section~\ref{sec:app})}. Recovering $\pi_{b[i]}$
from the response together with $\eseed_i$ is the column-extraction argument
of Lemma~1 in~\cite{FaultToForge2025}, structurally the analogue
of~\eqref{eq:meds-extract}.

\textbf{CROSS}~\cite{CROSSspec} builds on the Restricted Syndrome Decoding
Problem~\cite{BaldiBPSWW24} through a $5$-pass protocol. Its second challenge
is a binary weight-$w$ vector, so $\calH$ and $\calR$ have exactly the shape
they have in MEDS. Each leaf seed expands into $(\mathbf{e}'_i,
\mathbf{u}'_i)$, and a hidden round publishes $\mathbf{v}_i = \mathbf{e} \star
(\mathbf{e}'_i)^{-1}$, where $\mathbf{e}$ is the long-term R-SDP solution; an
attacker holding $\mathbf{v}_i$ and $\eseed_i$ recovers $\mathbf{e} =
\mathbf{v}_i \star \mathbf{e}'_i$ directly, which is how
ZKFault~\cite{ZKFault} achieves key recovery from one fault. CROSS also
authenticates the $w$ commitments the verifier cannot recompute with a Merkle
tree; this is orthogonal to seed hiding.

\textbf{MPCitH and VOLEitH schemes} use the tree in a superficially different
but structurally identical way. For each of $\tau$ repetitions the challenge
selects one hidden party index, and the signer publishes the \emph{co-path} of
siblings from that leaf to the root, letting the verifier recompute the other
$N-1$ shares. Mirath and RYDE use a forest of $\tau$ trees; FAEST and SDitH~v2
use a single tree with $\tau N$ leaves~\cite{OneTree}. The all-but-one vector
commitment is the counterpart of the flag tree, with $|\calH| = 1$ per
repetition. We include this family in the abstraction but do not attack it
experimentally.
